% --- Archon physics formalization source begin ---
% archon:physics
% archon:covers PhyXMiniProblems/problem_phyx_mini_0822.lean
% archon:source-report reports/phyx_mini/problem_phyx_mini_0822.source.json

\chapter{Physics problem phyx\_mini\_0822}
\label{ch:PhyXMiniProblems_problem_phyx_mini_0822}

\paragraph{Problem source.}
\#\# Physical scenario

Figure shows a hollow cylinder of uniform mass density \textbackslash{}( \textbackslash{}rho \textbackslash{}) with length \textbackslash{}( L \textbackslash{}), inner radius \textbackslash{}( R\_1 \textbackslash{}), and outer radius \textbackslash{}( R\_2 \textbackslash{}). (It might be a steel cylinder in a printing press.)

\#\# Question

Find its moment of inertia about its axis of symmetry

\#\# Answer choices

- A: A: \textbackslash{}(\textbackslash{}frac\{1\}\{3\}M(R\_1\textasciicircum{}2 + R\_2\textasciicircum{}2) \textbackslash{})

- B: B: \textbackslash{}(\textbackslash{}frac\{1\}\{2\}M(R\_1\textasciicircum{}2 + R\_2\textasciicircum{}2) \textbackslash{})

- C: C: \textbackslash{}(\textbackslash{}frac\{1\}\{4\}M(R\_1\textasciicircum{}2 + R\_2\textasciicircum{}2) \textbackslash{})

- D: D: \textbackslash{}(\textbackslash{}frac\{1\}\{5\}M(R\_1\textasciicircum{}2 + R\_2\textasciicircum{}2) \textbackslash{})

\#\# Figure caption (auxiliary; use the image as primary evidence)

The image is a 3D diagram of a cylindrical object with detailed annotations. It features:

1. **Cylinder Structure**:
   - An outer cylindrical shell with an inner cylindrical section.
   - The outer surface is shaded in yellow, while the inner section is shaded in pink.

2. **Axis**: 
   - A vertical dashed line, labeled "Axis," represents the central axis of the cylinder.

3. **Dimensions**:
   - The cylinder's height is labeled "L."
   - Two radial distances, \textbackslash{}( R\_1 \textbackslash{}) and \textbackslash{}( R\_2 \textbackslash{}), are marked from the axis to different points of the cylinder.

4. **Radial Segment**:
   - A small radial segment is marked with an arrow labeled \textbackslash{}( r \textbackslash{}) and a differential segment labeled \textbackslash{}( dr \textbackslash{}).

5. **Dashed Lines**:
   - Indicate boundaries and important features within the structure for clarity.

These elements suggest the diagram represents a cross-sectional view of a hollow cylindrical object, likely used in a physics or engineering context to illustrate concepts such as volume, surface area, or material distribution.

\#\# Dataset metadata

Rotational Motion; Physical Model Grounding Reasoning, Numerical Reasoning

\paragraph{Recorded answer/context.}
B: B: \textbackslash{}(\textbackslash{}frac\{1\}\{2\}M(R\_1\textasciicircum{}2 + R\_2\textasciicircum{}2) \textbackslash{})

\paragraph{Figure/image path.}
/root/proposal\_for\_physic/science-mango/phyx\_mini\_run/phyx\_data/test\_image/822.png

\paragraph{Formalization target.}
create a compiling Lean file with sorry bodies at `PhyXMiniProblems/problem\_phyx\_mini\_0822.lean`. The Lean declarations must preserve the physical quantities, dimensions or dimensional roles, figure labels, governing-law hypotheses, and final relation expressed by this problem.
Use Mathlib/Physlib names found through LeanExplore where available. If a physics API is missing, introduce faithful local abstractions rather than scalar placeholder aliases.

\begin{theorem}[Physics formalization target]
\label{thm:physics:phyx_mini_0822:target}
\lean{PhyXMiniProblems.ProblemPhyXMini0822.problem_phyx_mini_0822}
\uses{def:physics:phyx-mini-0822:phyxminiproblems-problemphyxmini0822-massinkilograms, def:physics:phyx-mini-0822:phyxminiproblems-problemphyxmini0822-lengthinmeters, def:physics:phyx-mini-0822:phyxminiproblems-problemphyxmini0822-momentofinertiainkilogrammeterssquared, def:physics:phyx-mini-0822:phyxminiproblems-problemphyxmini0822-uniformhollowcylindersetup, def:physics:phyx-mini-0822:phyxminiproblems-problemphyxmini0822-matchesproblemandfigure, def:physics:phyx-mini-0822:phyxminiproblems-problemphyxmini0822-hasphysicalcylinderparameters, def:physics:phyx-mini-0822:phyxminiproblems-problemphyxmini0822-satisfiesuniformannularcylinderlaws, def:physics:phyx-mini-0822:phyxminiproblems-problemphyxmini0822-matchesanswerchoice}
The assigned autoformalize agent should translate this physics problem into Lean declarations in the covered file, with theorem and lemma proof bodies written as `by sorry`.
\end{theorem}
\begin{proof}
This is an autoformalization task, not a proof task. Produce statements that can later be proved by the physics prover without weakening the physical model.
\end{proof}

% --- Archon declaration topology begin ---
\section{Declaration topology}
\begin{definition}[Mass Quantity]
  \label{def:physics:phyx-mini-0822:phyxminiproblems-problemphyxmini0822-massquantity}
  \lean{PhyXMiniProblems.ProblemPhyXMini0822.MassQuantity}
  A nonnegative physical mass
\end{definition}
\begin{proof}
  This is the declared model definition of Mass Quantity.
\end{proof}
\begin{definition}[Length Quantity]
  \label{def:physics:phyx-mini-0822:phyxminiproblems-problemphyxmini0822-lengthquantity}
  \lean{PhyXMiniProblems.ProblemPhyXMini0822.LengthQuantity}
  A nonnegative physical length, used for L, R₁, and R₂
\end{definition}
\begin{proof}
  This is the declared model definition of Length Quantity.
\end{proof}
\begin{definition}[Mass Density Quantity]
  \label{def:physics:phyx-mini-0822:phyxminiproblems-problemphyxmini0822-massdensityquantity}
  \lean{PhyXMiniProblems.ProblemPhyXMini0822.MassDensityQuantity}
  A nonnegative volume mass density, with dimension mass per length cubed
\end{definition}
\begin{proof}
  This is the declared model definition of Mass Density Quantity.
\end{proof}
\begin{definition}[Moment Of Inertia Quantity]
  \label{def:physics:phyx-mini-0822:phyxminiproblems-problemphyxmini0822-momentofinertiaquantity}
  \lean{PhyXMiniProblems.ProblemPhyXMini0822.MomentOfInertiaQuantity}
  A scalar moment of inertia, with dimension mass times length squared
\end{definition}
\begin{proof}
  This is the declared model definition of Moment Of Inertia Quantity.
\end{proof}
\begin{definition}[mass In Kilograms]
  \label{def:physics:phyx-mini-0822:phyxminiproblems-problemphyxmini0822-massinkilograms}
  \lean{PhyXMiniProblems.ProblemPhyXMini0822.massInKilograms}
  \uses{def:physics:phyx-mini-0822:phyxminiproblems-problemphyxmini0822-massquantity}
  Read a physical mass in coherent SI kilograms
\end{definition}
\begin{proof}
  This is the declared model definition of mass In Kilograms.
\end{proof}
\begin{definition}[length In Meters]
  \label{def:physics:phyx-mini-0822:phyxminiproblems-problemphyxmini0822-lengthinmeters}
  \lean{PhyXMiniProblems.ProblemPhyXMini0822.lengthInMeters}
  \uses{def:physics:phyx-mini-0822:phyxminiproblems-problemphyxmini0822-lengthquantity}
  Read a physical length in coherent SI meters
\end{definition}
\begin{proof}
  This is the declared model definition of length In Meters.
\end{proof}
\begin{definition}[mass Density In Kilograms Per Cubic Meter]
  \label{def:physics:phyx-mini-0822:phyxminiproblems-problemphyxmini0822-massdensityinkilogramspercubicmeter}
  \lean{PhyXMiniProblems.ProblemPhyXMini0822.massDensityInKilogramsPerCubicMeter}
  \uses{def:physics:phyx-mini-0822:phyxminiproblems-problemphyxmini0822-massdensityquantity}
  Read a volume mass density in coherent SI kilograms per cubic meter
\end{definition}
\begin{proof}
  This is the declared model definition of mass Density In Kilograms Per Cubic Meter.
\end{proof}
\begin{definition}[moment Of Inertia In Kilogram Meters Squared]
  \label{def:physics:phyx-mini-0822:phyxminiproblems-problemphyxmini0822-momentofinertiainkilogrammeterssquared}
  \lean{PhyXMiniProblems.ProblemPhyXMini0822.momentOfInertiaInKilogramMetersSquared}
  \uses{def:physics:phyx-mini-0822:phyxminiproblems-problemphyxmini0822-momentofinertiaquantity}
  Read a moment of inertia in coherent SI kilogram meters squared
\end{definition}
\begin{proof}
  This is the declared model definition of moment Of Inertia In Kilogram Meters Squared.
\end{proof}
\begin{definition}[Figure Feature]
  \label{def:physics:phyx-mini-0822:phyxminiproblems-problemphyxmini0822-figurefeature}
  \lean{PhyXMiniProblems.ProblemPhyXMini0822.FigureFeature}
  Physical features visibly distinguished in the supplied diagram
\end{definition}
\begin{proof}
  This is the declared model definition of Figure Feature.
\end{proof}
\begin{definition}[Figure Label]
  \label{def:physics:phyx-mini-0822:phyxminiproblems-problemphyxmini0822-figurelabel}
  \lean{PhyXMiniProblems.ProblemPhyXMini0822.FigureLabel}
  Text and mathematical labels printed in the supplied diagram
\end{definition}
\begin{proof}
  This is the declared model definition of Figure Label.
\end{proof}
\begin{definition}[Cylinder Direction]
  \label{def:physics:phyx-mini-0822:phyxminiproblems-problemphyxmini0822-cylinderdirection}
  \lean{PhyXMiniProblems.ProblemPhyXMini0822.CylinderDirection}
  The two geometrically relevant directions in the cylinder
\end{definition}
\begin{proof}
  This is the declared model definition of Cylinder Direction.
\end{proof}
\begin{definition}[Hollow Cylinder Figure]
  \label{def:physics:phyx-mini-0822:phyxminiproblems-problemphyxmini0822-hollowcylinderfigure}
  \lean{PhyXMiniProblems.ProblemPhyXMini0822.HollowCylinderFigure}
  \uses{def:physics:phyx-mini-0822:phyxminiproblems-problemphyxmini0822-figurefeature, def:physics:phyx-mini-0822:phyxminiproblems-problemphyxmini0822-figurelabel, def:physics:phyx-mini-0822:phyxminiproblems-problemphyxmini0822-cylinderdirection}
  Direct qualitative and label data transcribed from the primary image
\end{definition}
\begin{proof}
  This is the declared model definition of Hollow Cylinder Figure.
\end{proof}
\begin{definition}[Cylinder Density Model]
  \label{def:physics:phyx-mini-0822:phyxminiproblems-problemphyxmini0822-cylinderdensitymodel}
  \lean{PhyXMiniProblems.ProblemPhyXMini0822.CylinderDensityModel}
  Mass-distribution models distinguished by the problem statement
\end{definition}
\begin{proof}
  This is the declared model definition of Cylinder Density Model.
\end{proof}
\begin{definition}[Cylinder Geometry]
  \label{def:physics:phyx-mini-0822:phyxminiproblems-problemphyxmini0822-cylindergeometry}
  \lean{PhyXMiniProblems.ProblemPhyXMini0822.CylinderGeometry}
  Geometric models for the material region of the cylinder
\end{definition}
\begin{proof}
  This is the declared model definition of Cylinder Geometry.
\end{proof}
\begin{definition}[Uniform Hollow Cylinder Setup]
  \label{def:physics:phyx-mini-0822:phyxminiproblems-problemphyxmini0822-uniformhollowcylindersetup}
  \lean{PhyXMiniProblems.ProblemPhyXMini0822.UniformHollowCylinderSetup}
  \uses{def:physics:phyx-mini-0822:phyxminiproblems-problemphyxmini0822-massquantity, def:physics:phyx-mini-0822:phyxminiproblems-problemphyxmini0822-lengthquantity, def:physics:phyx-mini-0822:phyxminiproblems-problemphyxmini0822-massdensityquantity, def:physics:phyx-mini-0822:phyxminiproblems-problemphyxmini0822-momentofinertiaquantity, def:physics:phyx-mini-0822:phyxminiproblems-problemphyxmini0822-hollowcylinderfigure, def:physics:phyx-mini-0822:phyxminiproblems-problemphyxmini0822-cylinderdensitymodel, def:physics:phyx-mini-0822:phyxminiproblems-problemphyxmini0822-cylindergeometry}
  All dimensionful quantities are independent fields. In particular, axialMomentOfInertia is not defined from the answer formula. The scalar function is the SI readout of the mass in a coaxial shell per unit increment of its meter-valued radial coordinate
\end{definition}
\begin{proof}
  This is the declared model definition of Uniform Hollow Cylinder Setup.
\end{proof}
\begin{definition}[Matches Problem And Figure]
  \label{def:physics:phyx-mini-0822:phyxminiproblems-problemphyxmini0822-matchesproblemandfigure}
  \lean{PhyXMiniProblems.ProblemPhyXMini0822.MatchesProblemAndFigure}
  \uses{def:physics:phyx-mini-0822:phyxminiproblems-problemphyxmini0822-figurefeature, def:physics:phyx-mini-0822:phyxminiproblems-problemphyxmini0822-figurelabel, def:physics:phyx-mini-0822:phyxminiproblems-problemphyxmini0822-uniformhollowcylindersetup}
  Problem-text and primary-image facts. No field gives the cylinder's moment of inertia or selects an answer choice
\end{definition}
\begin{proof}
  This is the declared model definition of Matches Problem And Figure.
\end{proof}
\begin{definition}[Has Physical Cylinder Parameters]
  \label{def:physics:phyx-mini-0822:phyxminiproblems-problemphyxmini0822-hasphysicalcylinderparameters}
  \lean{PhyXMiniProblems.ProblemPhyXMini0822.HasPhysicalCylinderParameters}
  \uses{def:physics:phyx-mini-0822:phyxminiproblems-problemphyxmini0822-massinkilograms, def:physics:phyx-mini-0822:phyxminiproblems-problemphyxmini0822-lengthinmeters, def:physics:phyx-mini-0822:phyxminiproblems-problemphyxmini0822-massdensityinkilogramspercubicmeter, def:physics:phyx-mini-0822:phyxminiproblems-problemphyxmini0822-uniformhollowcylindersetup}
  Positivity and nondegeneracy of the independent cylinder parameters
\end{definition}
\begin{proof}
  This is the declared model definition of Has Physical Cylinder Parameters.
\end{proof}
\begin{definition}[Satisfies Uniform Annular Cylinder Laws]
  \label{def:physics:phyx-mini-0822:phyxminiproblems-problemphyxmini0822-satisfiesuniformannularcylinderlaws}
  \lean{PhyXMiniProblems.ProblemPhyXMini0822.SatisfiesUniformAnnularCylinderLaws}
  \uses{def:physics:phyx-mini-0822:phyxminiproblems-problemphyxmini0822-massinkilograms, def:physics:phyx-mini-0822:phyxminiproblems-problemphyxmini0822-lengthinmeters, def:physics:phyx-mini-0822:phyxminiproblems-problemphyxmini0822-massdensityinkilogramspercubicmeter, def:physics:phyx-mini-0822:phyxminiproblems-problemphyxmini0822-momentofinertiainkilogrammeterssquared, def:physics:phyx-mini-0822:phyxminiproblems-problemphyxmini0822-uniformhollowcylindersetup}
  General shell-integration laws for a uniform annular cylinder. At a radius r, the differential shell volume per radial increment is 2 π r L, so its mass per radial increment is ρ (2 π r L). Total mass is the radial integral of this shell profile, and axial moment of inertia is its radial second moment. The two bridge equalities connect these physical quantities to Physlib's RigidBody.mass and diagonal inertiaTensor entry.
\end{definition}
\begin{proof}
  This is the declared model definition of Satisfies Uniform Annular Cylinder Laws.
\end{proof}
\begin{definition}[Answer Choice]
  \label{def:physics:phyx-mini-0822:phyxminiproblems-problemphyxmini0822-answerchoice}
  \lean{PhyXMiniProblems.ProblemPhyXMini0822.AnswerChoice}
  Labels of the four alternatives printed with the problem
\end{definition}
\begin{proof}
  This is the declared model definition of Answer Choice.
\end{proof}
\begin{definition}[coefficient]
  \label{def:physics:phyx-mini-0822:phyxminiproblems-problemphyxmini0822-answerchoice-coefficient}
  \lean{PhyXMiniProblems.ProblemPhyXMini0822.AnswerChoice.coefficient}
  \uses{def:physics:phyx-mini-0822:phyxminiproblems-problemphyxmini0822-answerchoice}
  The coefficient multiplying M (R₁² + R₂²) in each printed choice
\end{definition}
\begin{proof}
  This is the declared model definition of coefficient.
\end{proof}
\begin{definition}[Matches Answer Choice]
  \label{def:physics:phyx-mini-0822:phyxminiproblems-problemphyxmini0822-matchesanswerchoice}
  \lean{PhyXMiniProblems.ProblemPhyXMini0822.MatchesAnswerChoice}
  \uses{def:physics:phyx-mini-0822:phyxminiproblems-problemphyxmini0822-massinkilograms, def:physics:phyx-mini-0822:phyxminiproblems-problemphyxmini0822-lengthinmeters, def:physics:phyx-mini-0822:phyxminiproblems-problemphyxmini0822-momentofinertiainkilogrammeterssquared, def:physics:phyx-mini-0822:phyxminiproblems-problemphyxmini0822-uniformhollowcylindersetup, def:physics:phyx-mini-0822:phyxminiproblems-problemphyxmini0822-answerchoice}
  A displayed choice agrees with the cylinder's independent inertia field
\end{definition}
\begin{proof}
  This is the declared model definition of Matches Answer Choice.
\end{proof}
% --- Archon declaration topology end ---
% --- Archon physics formalization source end ---
